
\section{Introduction}
Large language models (LLMs) are increasingly deployed as general-purpose orchestrators that coordinate tools and agents to solve complex tasks \citep{hong_metagpt_2024, tran_multi-agent_2025}.
A central pattern in this paradigm is collaboration with \emph{subagents} or agents that excel within a narrow domain ~\cite{anthropic_claude_subagents_2025,openai_codex_subagents_2025}.
This collaboration allows orchestrator LLMs to utilize the subagent's domain specific intelligence \citep{hong_metagpt_2024} and preserve its context length by task delegation \citep{zhang_chain_2024}, enabling effective multi-step reasoning and planning over long horizons.
Today, this collaboration is mediated entirely through natural language: the LLM invokes the subagent, receives a natural language description of its output, and reasons over that description to decide subsequent actions\citep{tran_multi-agent_2025}. 
Natural language collaboration is reasonable when both agents are themselves language models, as they share the same input space. 
However, in many important domains such as game playing, robotics, and  autonomous driving the strongest available agents are not language models; their expertise is encoded in internal representations: latent states capturing policy, value estimates, and learned features~\cite{silver2017mastering,monroe2024mastering,brohan2022rt}.
This mismatch between LLMs and non language subagents's input space raises a fundamental question: 
\begin{quote}
\centering\emph{How can LLMs effectively collaborate and jointly reason with non-language subagents?}
\end{quote}

As LLMs increasingly serve as general purpose orchestrators, the depth of their collaboration with these subagents determines what these systems can ultimately achieve.
Effective collaboration has the potential to unlock creative applications that neither achieves alone.
Consider chess: LLMs have been trained on more chess literature than most experts will study in a lifetime, yet they cannot leverage it to play the game competently, trailing far behind even amateur players~\cite{XXXChessBench}.
Conversely, pretrained engines surpassed human grandmasters decades ago~\cite{campbell2002deep}, yet they remain narrow specialists that are unable to explain the rationale behind a move or strategize under different contexts.
The most valued applications like commentating on a game, preparing against an opponent's known weaknesses, or making a creative puzzle therefore remain unsolved, because they demand both the engine's deep positional understanding and the LLM's ability to reason over human intent in natural language.
Improving how the LLM and the specialized subagent reason and collaborate would unlock these applications.

Existing approaches to LLM-subagent collaboration attempt to bridge this gap symbolically, by \emph{verbalizing} the subagent's outputs into natural language before passing them to the LLM.
Early work finetunes language models on textual descriptions of agent actions and value estimates \citep{zang2019automated, lee2022improving}, while more recent systems rely on in-context learning and tool-calling interfaces to surface subagent outputs at inference time \citep{schick2023toolformerlanguagemodelsteach, yao2022react, kim2024bridging}.
However, these approaches share a common assumption: that the subagent's expertise can be faithfully verbalized.
We argue that this assumption is fundamentally limiting.
A chess engine's latent representation of a state reflects millions of learned features like threats, tempo, pawn structure, and long-range tactical motifs as shown by \citet{jenner2024evidence} that cannot be losslessly compressed into a textual rationale.
Verbalization therefore forces the subagent's representations through a lossy bottleneck, collapsing rich latent structure into surface level descriptions.
Moreover, this error compounds across interactions, in multi-step settings, each exchange between the agents strips away details, and accumulates errors over the reasoning horizon, eroding the very benefits that motivated the collaboration in the first place.
Consequently, existing methods prevent orchestrator LLMs from fully exploiting the capabilities encoded in non-language subagents.

In this work we introduce the concept of \textit{latent state internalization} to bridge this gap.
Rather than verbalizing the domain agent's output, we inject its internal representations directly into the LLM's chain of thought as contiguous \textit{latent tokens}.
A lightweight projection module \textbf{LatentBridge}, projects these representations into the LLM's embedding space.
This enables the model to reason autoregressively over both agents' latent space simultaneously. \somesh{Open to alternate name suggestions for LatentBridge}

\ref{XXXfig1} describes the multi-step flow.
At each step the LLM (1) generates language to reason verbally, (2) updates the environment with actions, or (3) queries the domain agent.
At each query the domain agent processes the updated state and its latent state is projected as $K=32$ contiguous latent tokens.
These latent tokens are appended to the reasoning trace giving the LLM a growing record of the agent's evaluations as the multi step interaction progresses.
This creates a closed loop: the agent operates inside the generation process, updating its perspective after each action. We train on these interleaved traces of \lang{language}, \action{action}, and \state{latent tokens} with reinforcement learning, enabling the model to collaborate strategically with the domain agent.

We instantiate this framework as \textbf{LLAMIA} (\textbf{L}arge \textbf{L}anguage and \textbf{A}ction \textbf{M}odels with \textbf{I}nternal \textbf{A}gents), pairing [LLM backbone] with Leela Chess Zero's BT4 network (240M parameters), and evaluate on \textsc{LLAMIA-Bench}---our curated benchmark spanning [XXX] public benchmarks and we extend it to in-the-wild settings with [XXX] tasks.
This spans [XXX] tasks including behavior cloning across skill levels, behavioral stylometry, difficulty estimation, move and game commentary, and instruction-conditioned planning. We ablate our method against dedicated task experts, frontier reasoning models and RL trained baselines with access to the domain agent. LLAMIA matches or exceeds all task experts and baselines with a single model, with [XXX\%] gains on [XXX]. We observe that the margin grows for tasks requiring deeper collaboration and multi-step interactions like game level commentary and instruction-conditioned planning. Commentary produced through internalization shows highest overlaps with the lines discussed by grandmasters. While imitating gameplay we observe that LLAMIA even imitates tricks and blunders akin to humans under time pressure. RL with LLAMIA shows quicker and substantially higher performance gains. Our baselines and experiments conclusively prove the bottleneck of verbalization even when controlled for compute and data.

Beyond task performance, the internalized representations are themselves interpretable: a single direction in the latent space, identified via sparse autoencoders, recovers the effect of a natural-language instruction when injected without any language prompt.

% =============================================================================
% Why chess? (Testbed motivation)
% =============================================================================

Despite growing interest in latent agent communication~\cite{XXXlatentsurvey} coordination with non language agents remain unaddressed. We hypothesize that a central reason is a lack of environments and structured evaluations that offer the necessary combination: diverse collaborative tasks with verifiable metrics, large-scale public data, and pre-trained agents. Chess, with its longstanding history in AI, is therefore uniquely positioned to provide all three. Humans have collaborated with engines for decades across commentary, opponent preparation, puzzle curation, and difficulty assessment--- real world tasks with established evaluation protocols, documented at scale through public game databases\cite{Lichess,Caissabase,Gameknot} and annotated corpora~\cite{feng2024chessgpt}. We validate these findings primarily in chess, with a cross-domain pilot on poker providing initial evidence of transferability. \somesh{Struggling to position this in the paper, most likely instead of a separate defensive section, lets distribute it across paras}

% =============================================================================
% Summary of Contributions
% =============================================================================

\vspace{-1mm}
\subsection{Summary of Contributions}
\vspace{-1mm}

\noindent Our contributions address each stage of this gap: \somesh{A statement about broader literature and impact should be added here}

\begin{enumerate}[leftmargin=*, labelsep=0.5em, topsep=2pt, itemsep=2pt]
    \item \textbf{Latent state internalization (Section~\ref{sec:methodology}):} A novel paradigm that replaces verbal agent communication with direct latent representation injection, enabling end-to-end reinforcement learning over interleaved \lang{language}, \action{action}, and \state{latent tokens} Chain of Thoughts. The resulting framework, \textbf{LLAMIA}, achieves state-of-the-art performance across diverse tasks and baselines.

    \item \textbf{\textsc{LLAMIA-Bench} (Section~\ref{sec:experiments}):} To evaluate LLM and domain-agent collaboration, we curate [XXX] chess benchmarks, spanning [XXX] tasks spanning behavior cloning, stylometry, difficulty estimation, commentary, and instruction-conditioned planning.

    \item \textbf{Verbalization as the bottleneck (Section~\ref{sec:experiments}):} Controlled ablations show that verbalized communication is suboptimal both in terms of performance and compute efficiency, and the gap grows for tasks requiring deeper collaboration and multi-step interactions.

    \item \textbf{Interpretable latent collaboration (Section~\ref{sec:experiments}):} Sparse autoencoders recover latent directions whose injection replicates natural-language instructions without any language prompt---mechanistic evidence that internalization encodes task-relevant structure, not an opaque shortcut.
\end{enumerate}

\vspace{-2mm}
